\chapter{Nguồn dữ liệu}\label{ch:nguon}

\section{Tên bộ dữ liệu, nơi lấy, giấy phép, cách thu thập}\label{sec:nguon}
\begin{table}[H]
\centering
\caption{Thông tin nguồn của bộ dữ liệu}\label{tab:nguon}
\small
\begin{tabular}{@{}l L{10.2cm}@{}}
\toprule
Mục & Thông tin \\
\midrule
Tên bộ dữ liệu & \emph{Melbourne Housing Snapshot} \\
Nơi lấy & Kaggle, tài khoản đăng tải \texttt{dansbecker}; thông tin bộ dữ liệu tại \url{https://www.kaggle.com/datasets/dansbecker/melbourne-housing-snapshot} \\
Tệp sử dụng & \texttt{melb\_data.csv}, lưu trong \texttt{data/} \\
Giấy phép & CC BY-NC-SA 4.0 (Ghi công -- Phi thương mại -- Chia sẻ tương tự); cần giữ thông tin ghi công và tuân thủ điều kiện giấy phép khi chia sẻ hoặc biến đổi dữ liệu \\
Cách thu thập dữ liệu gốc & Bằng phương pháp cào dữ liệu web (web scraping) các tin đăng bán nhà công khai qua trang \url{https://www.domain.com.au}; bộ dữ liệu là bản chụp (snapshot) từ bộ Melbourne Housing của Tony Pino \\
Cách lấy dữ liệu trong dự án & Tải thủ công một tệp CSV từ Kaggle rồi đặt vào \texttt{data/}; dự án không tự cào dữ liệu hay gọi API để lấy bộ dữ liệu này. \\
\bottomrule
\end{tabular}
\end{table}

Mỗi dòng của tệp là một giao dịch bán nhà ở Melbourne (Úc), gồm vị trí, các thuộc tính của căn nhà, giá bán và ngày bán. Giá tính bằng đô la Úc (AUD).

Mọi số liệu trong báo cáo này chỉ tính trên tệp \texttt{melb\_data.csv}.

\section{Quy mô: số bản ghi, số thuộc tính, khoảng thời gian}\label{sec:quymo}
\begin{table}[H]
\centering
\caption{Quy mô và phạm vi của dữ liệu}\label{tab:quymo}
\small
\begin{tabular}{@{}l r@{}}
\toprule
Chỉ tiêu & Giá trị \\
\midrule
Số bản ghi (giao dịch) & 13,580 \\
Số thuộc tính (cột) trong tệp gốc & 21 \\
Số thuộc tính sau khi thêm biến dẫn xuất (\texttt{SaleMonth}, \texttt{HouseAge}) & 23 \\
Khoảng thời gian bán & 28/01/2016 -- 23/09/2017 \\
Số suburb & 314 \\
Số mã bưu chính (\texttt{Postcode}) & 198 \\
Số vùng (\texttt{Regionname}) & 8 \\
Số hội đồng quản lý (\texttt{CouncilArea}, chưa kể giá trị thiếu) & 33 \\
Số đại lý bán (\texttt{SellerG}) & 268 \\
Giá thấp nhất / cao nhất & 85,000 / 9,000,000 AUD \\
Giá trung vị / trung bình & 903,000 / 1,075,684 AUD \\
Số dòng trùng lặp hoàn toàn & 0 \\
\bottomrule
\end{tabular}
\end{table}

Bộ dữ liệu vượt xa quy mô tối thiểu mà đề tài yêu cầu (khoảng 500 bản ghi và 5 thuộc tính). Về cơ cấu, nhà riêng (\texttt{h}) chiếm 69.6\% số dòng (9,449 dòng), căn hộ (\texttt{u}) 22.2\% (3,017 dòng) và nhà liền kề (\texttt{t}) 8.2\% (1,114 dòng). Theo vùng, Southern Metropolitan chiếm 34.6\%, Northern Metropolitan 28.6\%, Western Metropolitan 21.7\% và Eastern Metropolitan 10.8\%; ba vùng Victoria (Eastern, Northern, Western) cộng lại chỉ có 126 dòng (0.9\%).

Về thời gian, dữ liệu trải trên 19 tháng có giao dịch nhưng không liên tục. Không có giao dịch nào trong tháng 3/2016 và tháng 1/2017; tháng 1/2016 chỉ có 2 giao dịch và tháng 2/2016 có 26 giao dịch, trong khi các tháng còn lại có từ 318 đến 1,536 giao dịch. Do đó quý 1/2016 gần như trống (28 giao dịch) và mọi so sánh theo thời gian cần đọc thận trọng (mục~\ref{sec:xuhuong}).

\section{Từ điển dữ liệu: ý nghĩa và kiểu từng cột}\label{sec:tudien}
Bảng~\ref{tab:tudien} mô tả 21 cột của tệp gốc. Ý nghĩa các cột theo cách hiểu thông dụng của bộ dữ liệu Melbourne Housing; đơn vị mét vuông của \texttt{Landsize} và \texttt{BuildingArea} được suy ra từ phân phối giá trị (trung vị sau làm sạch lần lượt là 533 và 126), vì tệp không kèm mô tả đơn vị. Tên hai cột \texttt{Lattitude} và \texttt{Longtitude} bị viết sai chính tả ngay từ tệp gốc và được giữ nguyên để không sửa dữ liệu thô.

{\small
\begin{longtable}{@{}l l L{6.6cm} L{3.1cm}@{}}
\caption{Từ điển dữ liệu của tệp \texttt{melb\_data.csv}}\label{tab:tudien}\\
\toprule
Cột & Kiểu & Ý nghĩa & Vai trò \\
\midrule
\endfirsthead
\toprule
Cột & Kiểu & Ý nghĩa & Vai trò \\
\midrule
\endhead
\midrule
\multicolumn{4}{r}{\emph{Tiếp ở trang sau}}\\
\endfoot
\bottomrule
\endlastfoot
\texttt{Suburb} & chuỗi & Tên suburb nơi căn nhà tọa lạc (314 giá trị) & Dùng, mã hóa mục tiêu \\
\texttt{Address} & chuỗi & Địa chỉ căn nhà; 13,378 giá trị khác nhau trên 13,580 dòng & Loại (gần như định danh) \\
\texttt{Rooms} & số nguyên & Số phòng (từ 1 đến 10) & Dùng \\
\texttt{Type} & chuỗi & Loại nhà: \texttt{h} nhà riêng, \texttt{t} nhà liền kề, \texttt{u} căn hộ (unit, duplex) & Dùng \\
\texttt{Price} & số thực & Giá bán, đơn vị AUD (85,000 đến 9,000,000) & Biến mục tiêu \\
\texttt{Method} & chuỗi & Phương thức bán: \texttt{S}, \texttt{SP}, \texttt{PI}, \texttt{VB}, \texttt{SA} & Loại (chỉ biết sau khi bán) \\
\texttt{SellerG} & chuỗi & Đại lý bán (268 giá trị) & Loại (chỉ biết sau khi bán) \\
\texttt{Date} & chuỗi & Ngày bán, định dạng \texttt{dd/mm/yyyy}, đọc thành kiểu ngày & Dùng để chia tập và tạo \texttt{SaleMonth} \\
\texttt{Distance} & số thực & Khoảng cách đến trung tâm thành phố, đơn vị km (0 đến 48.1) & Dùng \\
\texttt{Postcode} & số thực & Mã bưu chính (198 giá trị), thực chất là mã nhóm & Loại (trùng thông tin suburb) \\
\texttt{Bedroom2} & số thực & Số phòng ngủ từ một nguồn khác; trùng \texttt{Rooms} ở 95\% số dòng & Loại (dư thừa) \\
\texttt{Bathroom} & số thực & Số phòng tắm & Dùng \\
\texttt{Car} & số thực & Số chỗ đỗ xe & Dùng \\
\texttt{Landsize} & số thực & Diện tích đất, m\textsuperscript{2} & Dùng, \texttt{log1p} \\
\texttt{BuildingArea} & số thực & Diện tích xây dựng, m\textsuperscript{2} & Dùng, \texttt{log1p} \\
\texttt{YearBuilt} & số thực & Năm xây dựng & Dùng; tạo \texttt{HouseAge} \\
\texttt{CouncilArea} & chuỗi & Hội đồng địa phương quản lý khu vực (33 giá trị, 10.1\% thiếu) & Dùng, mã hóa mục tiêu \\
\texttt{Lattitude} & số thực & Vĩ độ & Dùng \\
\texttt{Longtitude} & số thực & Kinh độ & Dùng \\
\texttt{Regionname} & chuỗi & Vùng của thành phố (8 giá trị) & Dùng, one-hot \\
\texttt{Propertycount} & số thực & Số bất động sản trong suburb & Dùng \\
\end{longtable}
}

Kiểu trong bảng là kiểu Pandas khi đọc tệp: 8 cột chuỗi, 1 cột số nguyên và 12 cột số thực. Các cột số thực như \texttt{Postcode}, \texttt{Bedroom2}, \texttt{YearBuilt} mang bản chất số nguyên hoặc mã nhóm nhưng được lưu dạng số thực.

\section{Nguyên tắc giữ dữ liệu thô nguyên vẹn trong \texttt{data/}}\label{sec:nguyenven}
\begin{itemize}[nosep]
  \item Tệp \texttt{data/melb\_data.csv} giữ đúng như khi tải về; không sửa tay, không ghi đè. Mọi chương trình chỉ đọc tệp này.
  \item Mọi biến đổi (làm sạch, thêm biến, chia tập, mã hóa) thực hiện bằng mã trong \texttt{ml\_pipeline/features.py} và các bước \texttt{02\_prep.py}; kết quả trung gian (các tệp train, validation, test) được ghi riêng vào \texttt{ml\_pipeline/data/} và có thể tạo lại bất cứ lúc nào bằng lệnh chạy lại bước đó.
  \item Làm sạch không xóa dòng nào: các giá trị không hợp lệ được chuyển thành giá trị thiếu để bước tiền xử lý xử lý có kiểm soát (mục~\ref{sec:quytac}).
  \item Dữ liệu thô nằm ở \texttt{data/}, tách khỏi thư mục mã \texttt{ml\_pipeline/}, nên mã có thể thay đổi thoải mái mà không động tới bản gốc.
\end{itemize}
